\section{Deep Q-Network}
\label{sec:dqn}

Pierwszym zastosowanym algorytmem jest Deep Q-Network
\cite{mnih_dqn_2015,sutton_barto_rl}. Sieć neuronowa aproksymuje
funkcję wartości akcji $Q(s,a)$ i dla otrzymanego stanu zwraca
wartość dla każdej z sześciu możliwych decyzji.

Model przyjmuje reprezentację State A lub State B. Zastosowano dwie
warstwy ukryte po 256 neuronów z funkcją aktywacji ReLU oraz warstwę
wyjściową o sześciu neuronach:

\begin{equation}
    d_{\mathrm{state}}
    \rightarrow 256
    \rightarrow 256
    \rightarrow 6,
    \qquad
    d_{\mathrm{state}} \in \{373,393\}.
\end{equation}

Nie wszystkie akcje są legalne w każdej fazie rozdania. Przed
wyborem decyzji wartości odpowiadające akcjom niedozwolonym są
maskowane. To samo ograniczenie stosowane jest podczas wyznaczania
celu dla następnego stanu.

Podczas treningu wykorzystano strategię epsilon-greedy. Agent
początkowo często wybiera losowe legalne akcje, a wraz z kolejnymi
krokami prawdopodobieństwo eksploracji jest liniowo zmniejszane
z $1{,}0$ do $0{,}05$. Podczas ewaluacji eksploracja jest wyłączona
i wybierana jest akcja o największej przewidywanej wartości.

\subsection{Uczenie sieci Q}

Przejścia pomiędzy stanami zapisywane są w buforze experience replay.
Do aktualizacji sieci losowany jest minibatch wcześniejszych
doświadczeń. Oprócz aktualizowanej sieci Q utrzymywana jest osobna
sieć docelowa, której parametry są okresowo synchronizowane.

Dla przejścia nieterminalnego cel Bellmana ma postać

\begin{equation}
    y_t =
    r_t +
    \gamma
    \max_{a' \in A(s_{t+1})}
    Q_{\theta^-}(s_{t+1},a'),
\end{equation}

natomiast dla przejścia terminalnego

\begin{equation}
    y_t = r_t.
\end{equation}

Sieć uczona jest poprzez minimalizację funkcji Smooth L1 pomiędzy
$Q_\theta(s_t,a_t)$ a wartością docelową $y_t$. Do aktualizacji
parametrów zastosowano optymalizator Adam.

Podstawowa konfiguracja wykorzystuje learning rate $0{,}001$,
$\gamma=0{,}99$, minibatch 64 przejść, bufor 10000 przejść oraz
synchronizację sieci docelowej co 1000 kroków. Przez pierwsze
1000 kroków bufor jest jedynie wypełniany.

Źródła losowości związane z inicjalizacją sieci, środowiskiem,
strategią eksploracji i próbkowaniem bufora wykorzystują osobne
generatory wyprowadzone z głównego seeda eksperymentu. Wyuczone
parametry sieci wraz z konfiguracją, reprezentacją stanu i wariantem
nagrody mogą zostać zapisane do checkpointu i później wykorzystane
do deterministycznej ewaluacji.